% Generated draft: requires mathematical and source review.
\documentclass[11pt]{article}
\usepackage[margin=25mm]{geometry}
\usepackage{fontspec}
\setmainfont{Latin Modern Roman}
\usepackage{amsmath,amssymb,mathtools,unicode-math}
\setmathfont{Latin Modern Math}
\usepackage{enumitem,needspace,xurl,hyperref}
\setlength{\parindent}{0pt}
\setlength{\parskip}{5pt}
\setlength{\emergencystretch}{3em}
\newcommand{\N}{\mathbb{N}}
\newcommand{\Z}{\mathbb{Z}}
\newcommand{\Q}{\mathbb{Q}}
\newcommand{\R}{\mathbb{R}}
\newcommand{\C}{\mathbb{C}}
\newcommand{\F}{\mathbb{F}}
\newcommand{\A}{\mathbb{A}}
\newcommand{\PP}{\mathbb{P}}
\newcommand{\E}{\mathbb{E}}
\newcommand{\eps}{\varepsilon}
\newcommand{\dd}{\,\mathrm{d}}
\begin{document}
\section*{UnsolvedMath ID: 800542}
\textbf{Source problem code:} catalogue entry
\subsection{Definitions and mathematical statement}
The supplied source asks for recognition and characterization of segment visibility graphs, but does not supply its definition of visibility. The following formulation is conditional on the explicit conventions below.

Let \(S=\{s_1,\ldots,s_n\}\) be a finite family of nondegenerate closed straight line segments in \(\R^2\), with \(s_i\cap s_j=\varnothing\) whenever \(i\ne j\). For distinct points \(p,q\in\R^2\), write
\[
(p,q)=\{(1-t)p+tq:0<t<1\}.
\]
Under the convention that visibility may be witnessed by arbitrary points of the two segments and that the connecting open segment must avoid every segment of \(S\), define the segment visibility graph \(\operatorname{Vis}(S)\) to have vertex set \(\{1,\ldots,n\}\) and adjacency relation
\[
\{i,j\}\in E(\operatorname{Vis}(S))
\quad\Longleftrightarrow\quad
\exists p\in s_i\ \exists q\in s_j:
(p,q)\cap\bigcup_{k=1}^{n}s_k=\varnothing,
\qquad i\ne j.
\]
\textbf{Recognition problem.} Given a finite simple undirected graph \(G\) by its adjacency matrix, determine whether there exists such a family \(S\) for which
\[
G\cong\operatorname{Vis}(S).
\]
\textbf{Characterization problem.} Give necessary and sufficient conditions on a finite simple undirected graph \(G\) for the existence of such a representation.

\textbf{Unresolved source qualifications.} The supplied excerpt does not confirm the conventions concerning closed nondegenerate segments, allowable witness points, or whether obstruction is tested against all segments or only segments other than the two being joined. These conventions require confirmation against the cited source before this conditional formulation can be treated as the intended statement. The excerpt specifies no running-time bound or particular form of characterization.
\subsection{Short English statement}
Given a graph by its adjacency matrix, decide whether its vertices can be represented by pairwise disjoint plane segments so that adjacency corresponds exactly to visibility. Also characterize the graphs admitting such representations. The precise visibility and segment conventions require source confirmation.
\subsection{Sources}
\begin{itemize}
\item Local UnsolvedMath registry: permanent ID \(800542\), source problem code \texttt{catalogue entry}, title ``Visibility Graphs of Points, Segments and Polygons --- Problem 12.''
\item Ghosh and Goswami, \emph{Unsolved Problems in Visibility Graphs of Points, Segments and Polygons} (2010), Problem 12. Supplied source URL: \url{https://arxiv.org/abs/1012.5187}. The registry attributes the problem to Ghosh and Goswami and gives the proposed year as 2010.
\item Supplied provenance records access on 2026-10-16 and extraction by \texttt{arxiv\_latex}. This is supplied provenance, not an independently verified source visit.
\item The supplied registry labels the problem ``open'' and includes an \texttt{unsolved} environment. These are supplied status labels, not an independently verified assessment of the literature.
\end{itemize}
\end{document}
